
\documentclass{article}
\usepackage{colortbl}
\usepackage{makecell}
\usepackage{multirow}
\usepackage{supertabular}

\begin{document}

\newcounter{utterance}

\centering \large Interaction Transcript for game `wordle\_withcritic', experiment `high\_frequency\_words\_clue\_with\_critic', episode 8 with Qwen3.5{-}4B.
\vspace{24pt}

{ \footnotesize  \setcounter{utterance}{1}
\setlength{\tabcolsep}{0pt}
\begin{supertabular}{c@{$\;$}|p{.15\linewidth}@{}p{.15\linewidth}p{.15\linewidth}p{.15\linewidth}p{.15\linewidth}p{.15\linewidth}}
    \# & $\;$A & \multicolumn{4}{c}{Game Master} & $\;\:$B\\
    \hline

    \theutterance \stepcounter{utterance}  
    & & \multicolumn{4}{p{0.6\linewidth}}{
        \cellcolor[rgb]{0.9,0.9,0.9}{
            \makecell[{{p{\linewidth}}}]{
                \texttt{\tiny{[P1$\langle$GM]}}
                \texttt{You are a language wizard who likes to guess words by using the given rules.} \\
\\ 
\texttt{Welcome to Wordle! You have six attempts to guess the target word, a valid English word of five lowercase letters (a{-}z). Please use the tags "explanation:" and "guess:" to provide a concise explanation for each guess.} \\
\\ 
\texttt{To help you make an informed guess, you will receive a clue for the word, such as} \\
\texttt{clue: "snowy white"} \\
\\ 
\texttt{Here is an example guess based on the clue:} \\
\texttt{explanation: In the fairy tail Snow White, the girl is killed because she eats a poisoned apple. And the word apple has 5 letters.} \\
\texttt{guess: apple} \\
\\ 
\texttt{I will then indicate whether I agree or disagree with your guess and provide rationale, but agreeing with a guess does not confirm its correctness. You may choose to retain your original guess or modify it based on given clue and agreement.} \\
\\ 
\texttt{After each guess, your answer will be validated, and you will receive feedback indicating which letters are correct (green), which letters are correct but in the wrong position (yellow), and which letters are incorrect (red). This feedback can be useful in determining which letters to include or exclude in your next guess.} \\
\\ 
\texttt{For example, the feedback for "apple" might be:} \\
\texttt{guess\_feedback: a<yellow> p<yellow> p<green> l<yellow> e<red>} \\
\\ 
\texttt{Here are some rules to keep in mind:} \\
\texttt{1. Please do not generate any guess feedback or clue on your own} \\
\texttt{2. Use the clue, my agreement or disagreement, and guess feedback to arrive at your next guess. This should be reflected in your explanation as well.} \\
\\ 
\texttt{Do not generate any other text.} \\
\\ 
\texttt{Let's begin with your first guess.} \\
\\ 
\texttt{clue: tv show since 1/14/52} \\
            }
        }
    }
    & \\ \\

    \theutterance \stepcounter{utterance}  
    & \multicolumn{4}{p{0.6\linewidth}}{
        \cellcolor[rgb]{0.95,0.95,0.95}{
            \makecell[{{p{\linewidth}}}]{
                \texttt{\tiny{[P2$\rightarrow$P1]}}
                \texttt{{[}DRAFT{]}} \\
\texttt{explanation: The clue "tv show since 1/14/52" refers to "I Love Lucy," which premiered on October 15, 1951, but is often associated with the early 50s. However, the specific date 1/14/52 might refer to a specific episode or a different show. Let's consider "The Andy Griffith Show" which started in 1960, "Bewitched" in 1964, "I Dream of Jeannie" in 1965. Wait, 1/14/52 is January 14, 1952. A popular show that started around that time is "The Adventures of Ozzie and Harriet" (1952). But let's try a more iconic show. Maybe "The Twilight Zone" (1959). Let's try "I Love Lucy" again, but the date doesn't match exactly. Perhaps the clue is referring to "The Honeymooners" (1955). Let's try a show that definitely started in 1952. "The Adventures of Superman" (1952). Let's guess "superman".} \\
\texttt{guess: superman} \\
            }
        }
    }
    & & \\ \\

    \theutterance \stepcounter{utterance}  
    & \multicolumn{4}{p{0.6\linewidth}}{
        \cellcolor[rgb]{0.95,0.95,0.95}{
            \makecell[{{p{\linewidth}}}]{
                \texttt{\tiny{[P2$\rightarrow$P1]}}
                \texttt{{[}SELF{-}FEEDBACK{]}} \\
\texttt{FORMAT: problem} \\
\texttt{RULES: problem} \\
\texttt{STRATEGY: problem} \\
\texttt{REASON: The answer uses incorrect tags ("explanation:" and "guess:" instead of "explanation:" and "guess:" as per the prompt's specific instruction to use those exact tags, though the tags themselves are correct, the word "superman" is not a valid 5{-}letter English word, violating the game's core constraint, and the strategy is poor because it guesses a proper noun that doesn't fit the word length requirement).} \\
\\ 
\texttt{Wait, looking closer at the prompt instructions: "Please use the tags "explanation:" and "guess:"". The draft uses "explanation:" and "guess:". So FORMAT is actually ok.} \\
\texttt{The word "superman" is 8 letters, not 5. This violates the rule "valid English word of five lowercase letters". So RULES is problem.} \\
\texttt{The strategy is bad because it guesses an invalid word.} \\
\texttt{Let's re{-}evaluate FORMAT. The prompt says: "Please use the tags "explanation:" and "guess:"". The draft has:} \\
\texttt{explanation: ...} \\
\texttt{guess: ...} \\
\texttt{This matches. So FORMAT is ok.} \\
\texttt{RULES: The word "superman" is 8 letters. The game requires 5 letters. So RULES is problem.} \\
\texttt{STRATEGY: Guessing an 8{-}letter word is a bad move. So STRATEGY is problem.} \\
\\ 
\texttt{Let's re{-}read the FORMAT requirement carefully. "Do not generate any other text." The draft has no extra text outside the tags. The tags are correct. So FORMAT is ok.} \\
\\ 
\texttt{Wait, I need to check if the draft output matches the required format exactly.} \\
\texttt{The prompt says: "Please use the tags "explanation:" and "guess:"".} \\
\texttt{The draft says:} \\
\texttt{explanation: The clue...} \\
\texttt{guess: superman} \\
\texttt{This looks correct for FORMAT.} \\
\\ 
\texttt{However, the word "superman" is 8 letters. The game requires 5 letters. This is a RULES violation.} \\
\texttt{The strategy is also poor because it fails to find a 5{-}letter word.} \\
\\ 
\texttt{So:} \\
\texttt{FORMAT: ok} \\
\texttt{RULES: problem} \\
\texttt{STRATEGY: problem} \\
\texttt{REASON: The guess "superman" is 8 letters long, violating the game's requirement for a 5{-}letter word.} \\
\\ 
\texttt{Wait, let me double check the FORMAT example.} \\
\texttt{Example 1: "I will open..." {-}> FORMAT: problem (because it didn't start with "> ").} \\
\texttt{Example 2: "CLUE: flocking..." {-}> FORMAT: ok.} \\
\texttt{The draft in the user prompt is:} \\
\texttt{explanation: The clue "tv show since 1/14/52" refers to "I Love Lucy," which premiered on October 15, 1951, but is often associated with the early 50s. However, the specific date 1/14/} \\
            }
        }
    }
    & & \\ \\

    \theutterance \stepcounter{utterance}  
    & \multicolumn{4}{p{0.6\linewidth}}{
        \cellcolor[rgb]{0.9,0.9,0.9}{
            \makecell[{{p{\linewidth}}}]{
                \texttt{\tiny{[P1$\rangle$GM]}}
                \texttt{FORMAT: ok} \\
\texttt{RULES: problem} \\
\texttt{STRATEGY: problem} \\
\texttt{REASON: The guess "superman" is 8 letters long, violating the game's requirement for a 5{-}letter word.} \\
            }
        }
    }
    & & \\ \\

    \theutterance \stepcounter{utterance}  
    & & & \multicolumn{2}{p{0.3\linewidth}}{
        \cellcolor[rgb]{0.9,0.9,0.9}{
            \makecell[{{p{\linewidth}}}]{
                \texttt{\tiny{[GM$|$GM]}}
                \texttt{Error: The response should always start with the keyword 'explanation:'} \\
            }
        }
    }
    & & \\ \\

    \theutterance \stepcounter{utterance}  
    & & & \multicolumn{2}{p{0.3\linewidth}}{
        \cellcolor[rgb]{0.9,0.9,0.9}{
            \makecell[{{p{\linewidth}}}]{
                \texttt{\tiny{[GM$|$GM]}}
                \texttt{game\_result = ABORT} \\
            }
        }
    }
    & & \\ \\

\end{supertabular}
}

\end{document}
